\section*{Chimie (7 points)}
En chimie, les transformations peuvent être étudiées en se basant sur un suivi temporel, une mesure physique ou en utilisant des techniques appropriées. Les résultats et les conclusions obtenus dépendent de la nature des couples mis en jeu.\\
Cet exercice vise :

\begin{itemize}
  \item le suivi temporel d'une transformation chimique ;
  \item l'étude de la réaction entre l'acide monochloroéthanoïque et l'eau ;
  \item la détermination de la concentration d'acide monochloroéthanoïque dans un médicament.
\end{itemize}

\section*{Partie 1 : Cinétique chimique}
On mélange initialement à $t_{0}=0$, une solution aqueuse de sulfate de fer(II)
(\ce{Fe$^{2+}$_{(aq)} + SO$_4^{2-}$_{(aq)}}) avec une solution aqueuse de
sulfate de mercure(II) (\ce{Hg$^{2+}$_{(aq)} + SO$_4^{2-}$_{(aq)}}).  Le volume
du mélange est $V=\qty{100}{\mL}$.

La transformation chimique qui se produit est modélisée par l'équation chimique
suivante~:
\begin{center}
    \ce{2Fe$^{2+}$_{(aq)} + 2Hg$^{2+}$_{(aq)} -> 2Fe$^{3+}$_{(aq)} + Hg$_2^{2+}$_{(aq)}}
\end{center}

\begin{donnees}
    Concentrations molaires effectives des ions dans le mélange à l'état
    initial~: \\
    $\left[\ce{Fe$^{2+}$_{(aq)}}\right]_{0}=\qty{1.0e-2}{\mol\per\L}$~;
    $\left[\ce{Hg$^{2+}$_{(aq)}}\right]_{0}=\qty{2.0e-2}{\mol\per\L}$.
\end{donnees}
\begin{minipage}{0.55\linewidth}
    Le suivi de la transformation permet de tracer la courbe représentant
    l'évolution temporelle de l'avancement $x$ de la réaction (figure
    ci-contre).

    \begin{enumerate}
        \item Vérifier que la valeur de l'avancement maximal est
              $x_{\max}=\qty{5e-4}{\mol}$.
        \item Définir le temps de demi-réaction $t_{1/2}$ et déterminer sa valeur
              graphiquement.
        \item Déterminer, en ($\unit{\mol\per\L\per\h}$), la vitesse volumique de
              réaction à $t_{0}=0$.
        \item Comment évolue la vitesse volumique de réaction au cours de cette
              transformation~?
    \end{enumerate}
\end{minipage}
\hfill
\begin{minipage}{0.44\linewidth}
\begin{tikzpicture}[]
    \def\ta{20}
    \def\xf{0.5}
    %\pgfmathparse{ln(2)*\ta}\let\tm=\pgfmathresult
    \pgfmathparse{\xf/\ta}\let\p=\pgfmathresult
    \begin{axis}[
            cycle list name=my color list,
            width=8cm,height=7cm,
            xlabel=$t~(\unit{\h})$,ylabel=$x~(\unit{\mmol})$,
            xmin=0,xmax=59,
            ymin=0,ymax=0.54,
            %xtick={},
            %xticklabels={},
            %ytick={},
            %yticklabels={},
            variable=t,
            samples=200,
        ]
        \addplot[domain=0:60] {\xf*(1-exp(-\t/\ta))};
        \addplot[domain=0:21,dashed,thick] {\p*\t};
    \end{axis}
\end{tikzpicture}
\end{minipage}





Partie 2 : Acide monochloroéthanoïque\\
\includegraphics[width=\textwidth]{c8da97a1-45e2-47e7-8c6c-2b3cd6681557-2(1)}
En dermatologie, pour traiter les verrues (الاليل), on utilise un traitement physique ou un traitement chimique à l'aide d'un réactif corrosif.\\
L'acide monochloroéthanoïque $\mathrm{ClCH}_{2}-\mathrm{CO}_{2} \mathrm{H}$ est le principe actif de l'un des médicaments traitant I verrues. On utilisera la notation simplifiéc $A H$ pour cet acide et $A^{-}$pour sa base conjuguée.

\begin{enumerate}
  \item Étude de la réaction entre l'acide monochloroéthanoïque et l'eau
\end{enumerate}

On considère une solution aqueuse ( $S$ ) d'acide monochloroéthanoïque de volume $V=100 \mathrm{~mL}$ contenant initialement la quantité de matière $n_{1}(A H)=5.10^{-3} \mathrm{~mol}$. La mesure à $25^{\circ} \mathrm{C}$ de la conductivité $\sigma$ de la solution ( $S$ )donne $\sigma=0,3 S . m^{-1}$.

\section*{Données :}
\begin{itemize}
  \item On rappelle que, pour des solutions diluées, l'expression de la conductivité $\sigma$ en fonction des conductivités molaires ioniques $\lambda_{1}$ et des concentrations molaires effectives des espèces chimiques ioniques $X$, dissoutes s'écrit $\sigma=\sum \lambda_{\lambda}\left[X_{1}\right]$.
  \item Conductivités molaires ioniques à $25^{\circ} \mathrm{C}$ :
\end{itemize}

$$
\lambda_{1}=\lambda_{H_{3} O^{+}}=35.10^{-3} \mathrm{~S} . \mathrm{m}^{2} . \mathrm{mol}^{-1} \quad ; \quad \lambda_{2}=\lambda_{\lambda}=4.10^{-3} \mathrm{~S} . \mathrm{m}^{2} . \mathrm{mol}^{-1}
$$

\begin{itemize}
  \item L'effet des ions hydroxyde $\mathrm{HO}_{(\mathrm{aq})}^{-}$sur la conductivité de la solution est négligeable.
\end{itemize}

\subsection*{1.1. Donner la définition d'un acide selon Brönsted.}
1.2. L'acide monochloroéthanoïque réagit avec l'eau selon l'équation chimique :

$$
A H_{(a q)}+H_{2} O_{(1)} \rightleftharpoons A_{(a q)}^{-}+H_{3} O_{(a q)}^{+}
$$

1.2.a. Recopier, sur votre copie, le tableau d'avancement ci-après, et le compléter.

\begin{center}
\begin{tabular}{|l|l|l|l|l|l|}
\hline
\multicolumn{2}{|c|}{Équation chimique} & \multicolumn{4}{|c|}{$A H_{(m)}+H_{2} O_{(1)} \rightleftharpoons A_{(e p)}^{*}+H_{3} O_{(m)}^{*}$} \\
\hline
État du système & Avancement (mol) & \multicolumn{4}{|c|}{Quantités de matière (mol)} \\
\hline
Initial & $x=0$ & ......... & excès & ........... & ............ \\
\hline
Intermédiaire & $x$ & ......... & excès & ........... & ............ \\
\hline
Ėquilibre & $\boldsymbol{x}_{\boldsymbol{\text { dg }}}$ & ......... & excès & ........... & ........... \\
\hline
\end{tabular}
\end{center}

1.2.b. Montrer que l'avancement de la réaction à l'état d'équilibre est $x_{i q}=\frac{\sigma . V}{\lambda_{1}+\lambda_{2}}$. Calculer sa valeur.\\
1.2.c. Vérifier que la valeur du taux d'avancement final de la réaction est $\tau=0,154$.\\
1.3. Montrer que la constante d'acidité $K_{A}$ du couple $\left(A H_{(a) 1} / A_{(a-1)}^{-}\right)$peut s'exprimer par la relation :\\
$K_{A}=\frac{\tau^{2} \cdot n_{,}(A H)}{V \cdot(1-\tau)}$.\\
1.4. Calculer la valeur de $K_{d}$.\\
2. Détermination de la concentration d'acide monochloroéthanoïque dans un médicament

Un médicament de traitement des verrues se présente sous forme d'une solution aqueuse ( $S_{0}$ ) d'acide monochloroéthanoïque de concentration molaire $C_{0}$.\\
La solution $\left(S_{0}\right)$ étant concentrée. Pour déterminer la valeur de $C_{0}$, on prépare par dilution une solution aqueuse $\left(S_{1}\right)$ de concentration molaire $C_{1}=\frac{C_{0}}{10}$.\\
On réalise le titrage du volume $V_{1}=10 m L$ de la solution $\left(S_{1}\right)$ par unc solution aqueuse $\left(S_{B}\right)$ d'hydroxyde de sodium $N a_{(\mathrm{aq})}^{*}+H O_{(\mathrm{mq})}^{-}$de concentration molaire $C_{n}=0,10 \mathrm{~mol} . L^{-1}$. Le volume de la solution ( $S_{H}$ ) versé à l'équivalence est $V_{H, E}=21 \mathrm{~mL}$.\\
2.1. Écrire l'équation de la réaction chimique du dosage supposée totale.\\
2.2. Calculer la valeur de $C_{0}^{\prime}$.


